\section{Experimental Protocol}
\label{sec:protocol}

\subsection{Datasets and sample sizes}

We evaluate on three benchmarks spanning reasoning, knowledge, and summarization:
\begin{enumerate}
    \item \textbf{GSM8K}~\cite{Cobbe2021} test subset: 300 samples.
    \item \textbf{MMLU}~\cite{Hendrycks2020} subset: 500 samples (5 subjects $\times$ 100 each).
    \item \textbf{CNN/DailyMail}~\cite{Narayan2018} subset: 200 samples.
\end{enumerate}

\noindent\textbf{Sampling policy:} Fixed random seed $= 42$, stratified sampling where applicable, and a frozen sample-ID manifest generated before experimentation. Each output CSV records the seed and sample identifier, enabling paired analysis at the sample level.

\subsection{Prompt set}

We use the following fixed prompt templates for each task:

\smallskip\noindent\textbf{GSM8K:}
\begin{quote}\small
You are a careful math assistant. Solve the problem step by step and end with
Final Answer: \textless number\textgreater.\\
Question: \{question\}
\end{quote}

\noindent\textbf{MMLU:}
\begin{quote}\small
You are a knowledgeable assistant. Choose exactly one option.\\
Question: \{question\}\\
Options: A.~\{A\} B.~\{B\} C.~\{C\} D.~\{D\}\\
Answer:
\end{quote}

\noindent\textbf{CNN/DailyMail:}
\begin{quote}\small
Summarize the following article in 3--4 sentences, preserving key facts and entities.\\
Article: \{article\}\\
Summary:
\end{quote}
\paragraph{Objective and hypotheses}
This study quantifies practical speculative-decoding gains on consumer hardware
under fixed data and prompt controls.
We test four hypotheses:
\begin{itemize}
    \item \textbf{H1:} Same-family speculative decoding yields positive net speedup
        ($S>1$) over autoregressive decoding.
    \item \textbf{H2:} Increasing draft length from $k{=}4$ to $k{=}8$ improves
        speedup, while $k{=}16$ may reduce gains due to additional drafting
        and verification overhead.
    \item \textbf{H3:} Increasing $k$ tends to improve accepted block size but can
        reduce end-to-end speed once drafting/verification overhead dominates.
    \item \textbf{H4:} Deterministic decoding is more runtime-efficient and quality
        stable than stochastic decoding for the same draft and $k$.
\end{itemize}

\paragraph{Hardware and software controls}
We report two consumer-GPU run families:
an NVIDIA RTX4090 desktop profile and an RTX5090-laptop profile.
Both use the same Python runtime entry points, prompt templates,
tokenizer settings, and serial execution policy to minimize cross-run
contention. We keep software dependencies pinned within each run family and
preserve frozen manifests for paired comparison.

\begin{table}[t]
\centering
\caption{Hardware profile for the desktop RTX4090 run family.}
\label{tab:hw-profile}
\small
\begin{tabular}{ll}
\hline
\textbf{Component} & \textbf{Specification} \\
\hline
GPU & NVIDIA RTX4090 GPU \\
GPU memory & 24 GB GDDR6X \\
Memory bandwidth & 1008 GB/s \\
CUDA cores & 16384 \\
Architecture & Ada Lovelace \\
Power profile & TGP class 450~W \\
Software stack & PyTorch + Transformers (pinned env) \\
\hline
\end{tabular}
\end{table}

\begin{table}[t]
\centering
\caption{Hardware profile for the RTX5090-laptop run family.}
\label{tab:hw-profile-5090-laptop}
\small
\begin{tabular}{ll}
\hline
\textbf{Component} & \textbf{Specification} \\
\hline
GPU & NVIDIA GeForce RTX 5090 Laptop GPU \\
GPU memory & 24 GB GDDR7 \\
Memory bandwidth & 896 GB/s \\
CUDA cores & 10496 \\
Architecture & NVIDIA Blackwell \\
Power profile & Laptop TGP class up to 150~W (OEM-dependent) \\
Software stack & PyTorch + Transformers (pinned env) \\
\hline
\end{tabular}
\end{table}

\paragraph{Runtime setup}
We keep fixed step shapes, pre-allocate static draft/verify buffers, and warm
up the inference path before measurement. During decoding, each cycle updates
buffer contents under shape guards and falls back to eager execution when
required. This preserves output semantics and reduces avoidable runtime-path
variance in short speculative cycles.

\paragraph{Model pairing}
We evaluate two same-family pairings.
For Qwen2.5, the target is Qwen2.5-3B-Instruct with
Qwen2.5-0.5B-Instruct and Qwen2.5-1.5B-Instruct drafts on the RTX4090 run
family. For Qwen3, the target is Qwen3-4B-Instruct with a
Qwen3-0.6B-Instruct draft on RTX4090 and RTX5090-laptop run families.
For each family, fixed-policy experiments sweep $k\in\{4,8,16\}$ under
deterministic/stochastic regimes with regime-matched autoregressive baselines.
All configurations process fixed manifests to preserve paired comparability.

\paragraph{Task-entropy interpretation}
We treat the three tasks as distinct decoding-entropy regimes to interpret
acceptance behavior: GSM8K is a lower-entropy chain reasoning with concentrated
next-token mass; MMLU is medium-entropy constrained QA; CNN/DailyMail is
higher-entropy open-ended summarization. This framing is used only for
interpretation of acceptance and quality trends, not as an optimization
objective.

\paragraph{Decoding regimes and grid}
Two regimes are evaluated:
\begin{itemize}
    \item \textbf{Deterministic:} greedy decoding ($T{=}0.0$, top-$p{=}1.0$; sampling disabled).
    \item \textbf{Stochastic:} sampling with $T{=}0.7$, top-$p{=}0.9$, and no top-$k$ truncation.
\end{itemize}
For the primary fixed-policy matrix, we sweep $k\in\{4,8,16\}$ under both
regimes (6 configurations), plus regime-matched autoregressive baselines.

\paragraph{Endpoints}
The primary endpoint is paired wall-clock speedup
$S = L_{\text{autoregressive}}/L_{\text{spec}}$ as reported by the experiment summary
artifact for each draft--$k$--regime combination.
Secondary endpoints include acceptance rate $\alpha$, effective accepted block
size $B_{\text{eff}}$, mean per-sample latency, throughput (tokens/s),
task-level quality deltas, time to first token (TTFT), time per output token
(TPOT), and 99th-percentile (P99) latency. Statistical reporting uses
seed-level summaries (mean$\pm$std), paired
bootstrap confidence intervals where applicable, and paired t-tests with
Holm-Bonferroni correction for multi-metric adaptive-$k$ run-set comparisons.
